\section{Parte experimental}

\subsection{Metodología}

\textbf{Todo se mide; nada se estima.} Cada experimento es un guion que se ejecuta con un
comando, recibe sus parámetros por línea de comandos ---tamaños, número de consultas,
semilla--- y guarda sus resultados crudos en \texttt{benchmarks/results/}; de esos archivos
salen las gráficas y \emph{todas las cifras de esta sección}. Los tamaños son los del
enunciado: tablas de 1\,000, 10\,000 y 100\,000 registros. Las técnicas comparadas reciben
los mismos datos y las mismas consultas. Entorno: macOS, Python~3.12, páginas de 4~KB; para
el experimento espacial, PostgreSQL~17 con PostGIS~3.6.

\begin{lstlisting}
python -m benchmarks.storage_benchmark --sizes 1000 10000 100000 --queries 100
python -m benchmarks.index_benchmark   --sizes 1000 10000 100000 --queries 100
python -m benchmarks.spatial_benchmark --sizes 1000 10000 100000 --queries 100 \
       --points-file data/datasets/lugares_peru.csv
\end{lstlisting}

\begin{resumen}
\textbf{Cómo leer las gráficas y las tablas.}
\begin{itemize}[nosep, leftmargin=*]
  \item Cada gráfica tiene un panel por operación. El eje horizontal es el \textbf{tamaño de
        la tabla} (1\,000, 10\,000 y 100\,000 filas); el vertical, lo que se mide, con su
        unidad: \textbf{ms} son milisegundos y \textbf{KiB}, kibibytes (1\,024 bytes). Hay una
        línea por técnica, y \textbf{cuanto más abajo, mejor}.
  \item Los dos ejes son \textbf{logarítmicos}: cada marca vale diez veces la anterior. Una
        línea plana significa que el coste no depende del tamaño; una que sube en diagonal,
        que crece en proporción a él.
  \item El número escrito al final de cada línea es su valor con la tabla de 100\,000 filas.
        Las tablas traen los tres tamaños; en negrita va el mejor valor de cada columna.
\end{itemize}
\end{resumen}

\subsection{Experimento 1: heap file frente a archivo secuencial}

Sobre la misma tabla se mide cuánto tarda insertar todas las filas, cuánto tardan 100
búsquedas por clave primaria, cuánto ocupa el archivo y cuánto tarda reorganizarlo
(figura~\ref{fig:almacenamiento} y tabla~\ref{tab:almacenamiento}).

\grafica[0.78\textwidth]{almacenamiento.png}{fig:almacenamiento}{Heap file frente a archivo
secuencial. En la búsqueda se ve el cruce: el heap sube en diagonal y el secuencial es casi
plano.}

\begin{table}[H]
\centering
\caption{Heap file frente a archivo secuencial. Tiempos en milisegundos (ms) y espacio en KiB: \textbf{menos es mejor}, y en negrita va el mejor valor de cada tamaño.}\label{tab:almacenamiento}
\small
\begin{tabular}{@{}llrrr@{}}
\toprule
 & & \multicolumn{3}{c}{\textbf{Número de filas de la tabla}} \\
\cmidrule(lr){3-5}
\textbf{Operación} & \textbf{Técnica} & \textbf{1\,000} & \textbf{10\,000} & \textbf{100\,000} \\
\midrule
Insertar todas las filas & heap file & \textbf{7.66} & \textbf{71.8} & \textbf{744} \\
{\footnotesize\color{gris}tiempo total, en ms} & archivo secuencial & 25.3 & 265 & 3\,838 \\
\addlinespace[3pt]
100 búsquedas por clave & heap file & \textbf{1.08} & 10.6 & 97.2 \\
{\footnotesize\color{gris}tiempo total, en ms} & archivo secuencial & 4.62 & \textbf{4.54} & \textbf{5.74} \\
\addlinespace[3pt]
Espacio que ocupa el archivo & heap file & \textbf{64.0} & \textbf{576} & \textbf{5\,720} \\
{\footnotesize\color{gris}en KiB} & archivo secuencial & 100 & 964 & 9\,608 \\
\addlinespace[3pt]
Reorganizar el archivo {\footnotesize\color{gris}(ms)} & archivo secuencial & 2.63 & 19.3 & 194 \\
\bottomrule
\end{tabular}
\end{table}

\begin{itemize}[nosep, leftmargin=*]
  \item \textbf{El cruce está entre 1\,000 y 10\,000 filas.} Con 1\,000 registros, recorrer el
        heap entero (1.08 ms las 100 búsquedas) gana a la búsqueda binaria del secuencial
        (4.62 ms). A partir de ahí el heap crece de forma lineal y el secuencial casi no se
        mueve: con 100\,000 filas \textbf{el secuencial busca 17 veces más rápido}.
  \item \textbf{El orden se paga al escribir.} Insertar cuesta 5.2 veces más en el secuencial:
        cada inserción localiza su página y desplaza bytes, y cada tanto se dispara una
        reorganización.
  \item \textbf{El secuencial ocupa 1.7 veces más}, y es deliberado: el factor de llenado del
        80\,\% deja hueco para las inserciones futuras.
  \item \textbf{Cuándo usar cada uno:} heap si se escribe mucho y se consulta por índice;
        secuencial si se consulta por clave o por rango y las escrituras son la minoría.
\end{itemize}

\subsection{Experimento 2: B+ agrupado, B+ no agrupado y hash extendible}

Se mide cuánto tarda construir el índice, 100 búsquedas por igualdad, 100 consultas por
rango, recorrer la tabla entera en orden, y borrar y reinsertar el 10\,\% de las filas; y
cuánto ocupa (figura~\ref{fig:indices} y tabla~\ref{tab:indices}). Las búsquedas de los
índices no agrupados \textbf{incluyen el salto al heap}, que es el coste real de usarlos.

\grafica{indices.png}{fig:indices}{B+ agrupado, B+ no agrupado y hash extendible. En el
rango y en el recorrido ordenado falta la línea del hash: no puede hacer ninguna de las dos
cosas.}

\begin{table}[H]
\centering
\caption{Los tres índices. Tiempos en milisegundos (ms) y espacio en KiB: \textbf{menos es mejor}, y en negrita va el mejor valor de cada tamaño. \emph{No soportado}: la técnica no puede resolver esa operación.}\label{tab:indices}
\small
\begin{tabular}{@{}llrrr@{}}
\toprule
 & & \multicolumn{3}{c}{\textbf{Número de filas de la tabla}} \\
\cmidrule(lr){3-5}
\textbf{Operación} & \textbf{Técnica} & \textbf{1\,000} & \textbf{10\,000} & \textbf{100\,000} \\
\midrule
Construir el índice & B+ agrupado & 32.6 & 584 & 9\,259 \\
{\footnotesize\color{gris}tiempo total, en ms} & B+ no agrupado & 319 & 3\,790 & 49\,742 \\
 & hash extendible & \textbf{27.1} & \textbf{339} & \textbf{3\,464} \\
\addlinespace[3pt]
100 búsquedas por igualdad & B+ agrupado & \textbf{1.72} & \textbf{6.23} & \textbf{7.98} \\
{\footnotesize\color{gris}tiempo total, en ms} & B+ no agrupado & 23.9 & 24.2 & 33.6 \\
 & hash extendible & 12.5 & 9.59 & 11.1 \\
\addlinespace[3pt]
100 consultas por rango & B+ agrupado & \textbf{3.41} & \textbf{8.66} & \textbf{11.0} \\
{\footnotesize\color{gris}tiempo total, en ms} & B+ no agrupado & 34.3 & 36.3 & 53.5 \\
 & hash extendible & \multicolumn{3}{c}{\emph{no soportado}} \\
\addlinespace[3pt]
Recorrer toda la tabla en orden & B+ agrupado & \textbf{0.31} & \textbf{3.50} & \textbf{36.9} \\
{\footnotesize\color{gris}tiempo total, en ms} & B+ no agrupado & 2.13 & 22.8 & 390 \\
 & hash extendible & \multicolumn{3}{c}{\emph{no soportado}} \\
\addlinespace[3pt]
Borrar y reinsertar el 10\,\% & B+ agrupado & \textbf{7.61} & 174 & 2\,295 \\
{\footnotesize\color{gris}tiempo total, en ms} & B+ no agrupado & 79.5 & 911 & 9\,622 \\
 & hash extendible & 15.0 & \textbf{111} & \textbf{1\,464} \\
\addlinespace[3pt]
Espacio que ocupa el índice & B+ agrupado & 100 & 924 & 9\,312 \\
{\footnotesize\color{gris}en KiB} & B+ no agrupado & \textbf{24.0} & \textbf{220} & 2\,060 \\
 & hash extendible & 28.0 & 264 & \textbf{2\,056} \\
\bottomrule
\end{tabular}
\end{table}

\begin{itemize}[nosep, leftmargin=*]
  \item \textbf{El B+ agrupado gana en consultas.} Responde 100 igualdades en
        7.98 ms con 100\,000 filas porque la fila ya está en la hoja. El hash tarda
        11.1 ms pese a ser $O(1)$: toda la diferencia es el salto al heap. El B+ no agrupado
        (33.6 ms) suma $O(\log N)$ \emph{más} el salto.
  \item \textbf{El salto al heap domina cuando hay muchos resultados.} Recorrer en orden
        100\,000 filas tarda 36.9 ms en el agrupado y 390 ms en el no agrupado:
        11 veces más con la misma estructura. La diferencia son 100\,000 accesos sueltos.
  \item \textbf{El hash gana al escribir} ---construye 14 veces más rápido que el B+ no
        agrupado y es el más rápido borrando y reinsertando--- \textbf{pero no admite rangos ni
        orden}: no es un cero de rendimiento, es una limitación de la técnica.
  \item \textbf{El espacio va al revés.} El agrupado ocupa 4.5 veces más que los otros dos
        porque guarda las filas; por eso solo puede haber uno por tabla y tantos secundarios
        como haga falta.
  \item \textbf{Cuándo usar cada uno:} agrupado para la clave primaria; hash para columnas que
        solo se consultan por igualdad; B+ no agrupado cuando además hacen falta rangos.
\end{itemize}

\subsection{Experimento 3: búsqueda secuencial, R-Tree y GiST de PostgreSQL}

Las tres técnicas responden las mismas consultas sobre los mismos puntos: \textbf{lugares
reales de OpenStreetMap} del Perú, de los que se toman muestras de 1\,000, 10\,000 y
100\,000. Las consultas son las del enunciado: los lugares a menos de 1, 5 y 10~km de un
punto, y los 10, 50 y 100 más cercanos a él; cien de cada una, centradas junto a puntos de
la muestra. Las dos técnicas del motor se miden con el mismo SQL de punta a punta; GiST, con
PostGIS sobre una columna \texttt{geography}, \texttt{ST\_DWithin} y el operador
\texttt{<->}.

\paragraph{Antes de medir: \textquestiondown{}responden lo mismo?} Un índice que acelera cambiando la
respuesta no es una comparación. El experimento aborta si el R-Tree devuelve filas distintas
de las del recorrido, y anota cuántas consultas de PostGIS coinciden con el motor:
\textbf{las 1\,800 consultas coinciden, fila por fila}, y PostgreSQL confirma con
\texttt{EXPLAIN} que usa el índice en todas.

\grafica{espacial-radio.png}{fig:radio}{Consultas por radio. La línea del recorrido
secuencial es la misma en los tres paneles: no depende del radio.}

\begin{table}[H]
\centering
\caption{Consultas por radio: tiempo medio de \emph{una} consulta, en milisegundos (promedio de 100). Menos es mejor; en negrita, el mejor valor de cada tamaño. La última columna es cuántos lugares devuelve de media cada consulta sobre la tabla de 100\,000.}\label{tab:radio}
\small
\begin{tabular}{@{}llrrrr@{}}
\toprule
 & & \multicolumn{3}{c}{\textbf{Número de puntos de la tabla}} & \\
\cmidrule(lr){3-5}
\textbf{Consulta} & \textbf{Técnica} & \textbf{1\,000} & \textbf{10\,000} & \textbf{100\,000} & \textbf{Lugares devueltos} \\
\midrule
A menos de 1 km & búsqueda secuencial & 3.67 & 34.5 & 353 & 64.1 \\
{\footnotesize\color{gris}tiempo por consulta, en ms} & R-Tree & 0.25 & 0.40 & 1.07 & \\
 & GiST (PostgreSQL) & \textbf{0.051} & \textbf{0.045} & \textbf{0.25} & \\
\addlinespace[3pt]
A menos de 5 km & búsqueda secuencial & 3.60 & 34.4 & 357 & 544 \\
{\footnotesize\color{gris}tiempo por consulta, en ms} & R-Tree & 0.29 & 0.75 & 5.49 & \\
 & GiST (PostgreSQL) & \textbf{0.045} & \textbf{0.061} & \textbf{0.50} & \\
\addlinespace[3pt]
A menos de 10 km & búsqueda secuencial & 3.63 & 34.9 & 353 & 1\,577 \\
{\footnotesize\color{gris}tiempo por consulta, en ms} & R-Tree & 0.38 & 1.30 & 11.5 & \\
 & GiST (PostgreSQL) & \textbf{0.051} & \textbf{0.10} & \textbf{1.00} & \\
\bottomrule
\end{tabular}
\end{table}

\grafica{espacial-knn.png}{fig:knn}{Vecinos más cercanos. El R-Tree y GiST son casi planos:
su coste apenas depende del tamaño de la tabla.}

\begin{table}[H]
\centering
\caption{Vecinos más cercanos: tiempo medio de \emph{una} consulta, en milisegundos (promedio de 100). Menos es mejor; en negrita, el mejor valor de cada tamaño.}\label{tab:knn}
\small
\begin{tabular}{@{}llrrr@{}}
\toprule
 & & \multicolumn{3}{c}{\textbf{Número de puntos de la tabla}} \\
\cmidrule(lr){3-5}
\textbf{Consulta} & \textbf{Técnica} & \textbf{1\,000} & \textbf{10\,000} & \textbf{100\,000} \\
\midrule
Los 10 más cercanos & búsqueda secuencial & 7.75 & 75.6 & 1\,302 \\
{\footnotesize\color{gris}tiempo por consulta, en ms} & R-Tree & 0.31 & 0.42 & 0.64 \\
 & GiST (PostgreSQL) & \textbf{0.058} & \textbf{0.056} & \textbf{0.13} \\
\addlinespace[3pt]
Los 50 más cercanos & búsqueda secuencial & 8.02 & 76.2 & 1\,304 \\
{\footnotesize\color{gris}tiempo por consulta, en ms} & R-Tree & 0.49 & 0.65 & 1.02 \\
 & GiST (PostgreSQL) & \textbf{0.10} & \textbf{0.096} & \textbf{0.21} \\
\addlinespace[3pt]
Los 100 más cercanos & búsqueda secuencial & 8.14 & 76.1 & 1\,308 \\
{\footnotesize\color{gris}tiempo por consulta, en ms} & R-Tree & 0.73 & 0.97 & 1.37 \\
 & GiST (PostgreSQL) & \textbf{0.12} & \textbf{0.13} & \textbf{0.28} \\
\bottomrule
\end{tabular}
\end{table}

\paragraph{Por qué: los nodos que abre cada consulta.} El plan de ejecución dice cuántos
nodos del R-Tree visitó cada búsqueda (figura~\ref{fig:nodos} y tabla~\ref{tab:nodos}). El árbol pasa de 8 a
610 nodos, y buscar los 10 más cercanos pasa de abrir 2.3 nodos a abrir 4.1: por eso
el tiempo casi no crece.

\grafica[0.74\textwidth]{espacial-nodos.png}{fig:nodos}{Nodos del R-Tree que abre de media cada
consulta, frente a todos los que tiene el árbol (línea gris discontinua).}

\begin{table}[H]
\centering
\caption{Nodos del R-Tree que abre de media cada consulta (promedio de 100), y nodos que tiene el
árbol. Cuantos menos abre, menos páginas lee de disco.}\label{tab:nodos}
\small
\begin{tabular}{@{}lrrr@{}}
\toprule
 & \multicolumn{3}{c}{\textbf{Número de puntos de la tabla}} \\
\cmidrule(l){2-4}
\textbf{Consulta} & \textbf{1\,000} & \textbf{10\,000} & \textbf{100\,000} \\
\midrule
    A menos de 1 km & 2.0 & 2.2 & 4.5 \\
    A menos de 5 km & 2.1 & 3.0 & 9.6 \\
    A menos de 10 km & 2.2 & 3.8 & 17.4 \\
    Los 10 más cercanos & 2.3 & 2.5 & 4.1 \\
    Los 50 más cercanos & 2.8 & 3.2 & 5.0 \\
    Los 100 más cercanos & 3.4 & 4.3 & 5.9 \\
\midrule
Nodos que tiene el árbol & 8 & 62 & 610 \\
\bottomrule
\end{tabular}
\end{table}

\grafica[0.90\textwidth]{espacial-indice.png}{fig:indice}{Construir el índice, guardarlo y
consultarlo. La línea discontinua es el mismo R-Tree construido insertando los puntos de
uno en uno en vez de con carga masiva.}

\begin{table}[H]
\centering
\caption{Lo que cuesta tener el índice espacial. Tiempos en milisegundos (ms), espacio y memoria en KiB: \textbf{menos es mejor}, y en negrita va el mejor valor de cada tamaño.}\label{tab:indice}
\small
\begin{tabular}{@{}llrrr@{}}
\toprule
 & & \multicolumn{3}{c}{\textbf{Número de puntos de la tabla}} \\
\cmidrule(lr){3-5}
\textbf{Medida} & \textbf{Técnica} & \textbf{1\,000} & \textbf{10\,000} & \textbf{100\,000} \\
\midrule
Construir el índice & R-Tree, carga masiva & 5.52 & 46.1 & 687 \\
{\footnotesize\color{gris}tiempo total, en ms} & R-Tree, inserción uno a uno & 132 & 1\,746 & 20\,922 \\
 & GiST (PostgreSQL) & \textbf{1.78} & \textbf{16.7} & \textbf{199} \\
\addlinespace[3pt]
Espacio que ocupa el índice & R-Tree, carga masiva & \textbf{36.0} & \textbf{252} & \textbf{2\,444} \\
{\footnotesize\color{gris}en KiB} & R-Tree, inserción uno a uno & 40.0 & 340 & 3\,264 \\
 & GiST (PostgreSQL) & 88.0 & 752 & 7\,664 \\
\addlinespace[3pt]
Memoria máxima al consultar & búsqueda secuencial & 200 & 592 & \textbf{1\,837} \\
{\footnotesize\color{gris}en KiB} & R-Tree & \textbf{94.4} & \textbf{166} & 2\,204 \\
\bottomrule
\end{tabular}
\end{table}

\paragraph{Conclusiones del experimento espacial.}
\begin{itemize}[nosep, leftmargin=*]
  \item \textbf{El recorrido crece con la tabla; el índice, con la respuesta.} La búsqueda
        secuencial tarda lo mismo pida 1~km o 10 (353 y 353 ms con 100\,000
        puntos), porque siempre mide la distancia a todos. El R-Tree tarda 1.07 ms para
        1~km y 11.5 ms para 10: lo que le cuesta son los lugares que devuelve.
  \item \textbf{Cuanto más selectiva la consulta, más gana el índice:} 328 veces más
        rápido que recorrer para 1~km, 65 para 5~km y 31 para 10~km, donde cada
        consulta devuelve ya el 1.6\,\% de la tabla.
  \item \textbf{En los vecinos más cercanos la diferencia es de tres órdenes de magnitud.} Sin
        índice hay que medir los 100\,000 puntos y ordenarlos: 1\,302 ms. El R-Tree abre
        los nodos del más cercano al más lejano y se detiene al tener diez: 0.64 ms,
        \textbf{2\,022 veces menos}. Y casi no depende del tamaño: de 1\,000 a 100\,000
        puntos el recorrido se multiplica por 168 y el R-Tree por 2.1.
  \item \textbf{GiST es entre 4.4 y 13 veces más rápido que el R-Tree, y escala igual.}
        Las dos líneas son paralelas en las gráficas. GiST está escrito en C y lleva décadas
        de optimización; este R-Tree es Python y deserializa cada nodo que abre. Lo que el
        experimento compara es \emph{el algoritmo}, y ahí las dos estructuras se comportan
        igual y el recorrido no.
  \item \textbf{La carga masiva es 30 veces más rápida que insertar y deja un árbol
        mejor:} 610 nodos en vez de 815 y un 25\,\% menos de espacio. El índice
        propio ocupa 3.1 veces menos que el de PostGIS, que indexa cajas en tres dimensiones.
\end{itemize}

\subsection{Conclusiones generales}

\begin{itemize}[nosep, leftmargin=*]
  \item \textbf{No hay una estructura mejor: hay una mejor para cada consulta.} El heap gana
        al insertar y el secuencial al buscar; el B+ agrupado gana en consultas y el hash en
        escrituras; el R-Tree gana tanto más cuanto más selectiva es la consulta (328 veces
        con 1~km, 31 con 10~km).
  \item \textbf{El coste que domina es el acceso a disco, y el que se olvida es el salto al
        heap:} la misma estructura, el B+, tarda 11 veces más en un recorrido ordenado
        cuando sus hojas apuntan a las filas en vez de contenerlas.
  \item \textbf{Un índice vale lo que deja de leer.} El R-Tree encuentra los 10 lugares más
        cercanos entre 100\,000 abriendo de media 4.1 de sus 610 nodos.
  \item \textbf{Medir cambia las conclusiones.} Sin los experimentos no se habría visto que
        el heap gana al secuencial con 1\,000 filas, que el hash pierde frente al B+ agrupado
        en igualdad por el salto al heap, ni cuánto mejora el árbol la carga masiva.
\end{itemize}
